% !TeX root = ../main.tex
\section{Краткое функциональное и математическое описание}

Рассматриваемое устройство осуществляет сквозную цифровую обработку радиолокационного сигнала, включающую аналого-цифровую субдискретизацию, цифровой перенос частоты, фильтрацию, децимацию и оптимальную (согласованную) фильтрацию.

\subsection{Аналого-цифровое преобразование и субдискретизация}
Входной аналоговый сигнал представляет собой радиоимпульс с линейной частотной модуляцией (ЛЧМ), расположенный на промежуточной частоте $f_c = 170$ МГц. Аналого-цифровой преобразователь (АЦП) осуществляет дискретизацию сигнала с частотой $f_s = 140$ МГц и квантование с разрядностью $n_a = 9$ бит. Поскольку частота дискретизации меньше удвоенной максимальной частоты спектра сигнала, происходит целенаправленное наложение спектров (алиасинг). Центральная частота оцифрованного сигнала смещается на разностную частоту:
$$ f_{alias} = |f_c - f_s| = 30 \text{ МГц}. $$
Математическая модель квантованного сигнала на выходе АЦП имеет вид:
$$ x_q[n] = Q_{n_a} \{ x(n T_s) \}, $$
где $T_s = 1/f_s$ --- период дискретизации, $Q_{n_a}$ --- оператор квантования для $9$ бит.

\subsection{Цифровая демодуляция (ЦДМ)}
Для переноса спектра сигнала с алиасной частоты на нулевую (в область видеочастот) применяется цифровой гетеродин (ЦГ) и комплексный умножитель. Генератор с числовым программным управлением (NCO) формирует комплексную экспоненту с центральной частотой $f_G = 170$ МГц и разрядностью $n_g = 9$ бит:
$$ s_{NCO}[n] = Q_{n_g} \{ \exp(-j 2\pi f_G n T_s) \}. $$
В дискретной области частота $170$ МГц при $f_s = 140$ МГц эквивалентна частоте $-30$ МГц, что обеспечивает точный перенос спектра полезного сигнала на $0$ Гц. Сигнал на выходе ЦДМ:
$$ y_{cdm}[n] = x_q[n] \cdot s_{NCO}[n]. $$

\subsection{Цифровая фильтрация (ЦФ) и децимация}
После перемножения образуются высокочастотные комбинационные составляющие. Для их подавления используется цифровой фильтр нижних частот (ФНЧ) с частотой пропускания, согласованной с половиной полосы сигнала $\Delta f / 2 = 10$ МГц.
$$ y_{f}[n] = \sum_{k=0}^{K-1} b_k \cdot y_{cdm}[n-k], $$
где $b_k$ --- коэффициенты импульсной характеристики ФНЧ.
Далее производится понижение частоты дискретизации (децимация) в $D$ раз для оптимизации вычислительной нагрузки в последующих трактах:
$$ y_{dec}[m] = y_{f}[m \cdot D]. $$

\subsection{Согласованная фильтрация (СФ)}
Оптимальная обработка принятого сигнала осуществляется в согласованном фильтре, импульсная характеристика которого $h[m]$ является зеркальной и комплексно-сопряженной копией излученного зондирующего ЛЧМ-сигнала.Для снижения уровня боковых лепестков сжатого импульса применяется весовая функция Хэмминга:
$$ h_w[m] = h[m] \cdot \left(0.54 - 0.46 \cos\left(\frac{2\pi m}{M-1}\right)\right), $$
где $M$ --- длина импульсной характеристики.

Для вычисления линейной свертки длинных последовательностей применяется вычислительно эффективный алгоритм перекрытия со сложением (Overlap-Add), реализуемый встроенной функцией \texttt{fftfilt}. Непрерывный входной поток разбивается на неперекрывающиеся блоки длиной $L$. Вычисляется $N$-точечное дискретное преобразование Фурье (ДПФ) каждого блока и импульсной характеристики (где $N = L + M - 1$). Векторы перемножаются в частотной области, после чего выполняется обратное ДПФ.
$$ Y_i[k] = \text{FFT}\{y_{dec\_i}[m]\} \cdot \text{FFT}\{h_w[m]\} $$
$$ y_{sf\_i}[m] = \text{IFFT}\{Y_i[k]\} $$
Перекрывающиеся "хвосты" соседних обработанных блоков (длиной $M-1$ отсчетов) складываются, формируя итоговый выходной сигнал без краевых искажений.